% v12 front matter and introduction. Every number is a generated macro from the v11 evidence build.
\newcommand{\vTwelveAbstract}{%
Review comments must follow revised code, and a relocated comment must not be mistaken for a settled concern. Review tools drop a changed line's position or move the comment without saying whether the reviewed code changed. We separate three properties that tools merge: whether the reviewed code changed, whether the anchor still points at it, and whether the supervisor's concern is resolved. We state five safety requirements over them, and a counterexample for each failure mechanism. A safety layer wraps any code locator: change labels computed from the reviewed quotation, indentation and copy guards on the returned window, and issue states that relocation cannot close. We tested it over a capstone review portal's matcher, the region mapper CodeMapper, a line tracker and two language models, on controlled edits, real and student-authored Python histories, a human-annotated benchmark, and \wkUnits{} GitHub review suggestions in five languages, whose application gives an annotator-free reference. CodeMapper located code better than the portal's matcher. Over CodeMapper the layer cut missed re-checks on real histories (syntax-tree reference) from \wiRealCmQMissK{} to \wiRealCmGiMissK{} of \wiRealCmQMissN{}, at the cost of confirmations on unchanged code. No locator showed a suggestion's target on wrong code without a request. Edits wrapping unchanged quoted lines evade every quotation check, a copy guard cannot tell identical copies apart, tuned thresholds did not transfer between projects, and the portal's lifecycle loses issues through deletion and dismissal; a model-checked contract closes these gaps. With no supervisor or student observed, the results show technical feasibility, not educational benefit.}

\begin{document}

\title{Living Review: Carrying Code Review Comments Safely Across Resubmissions}
\author{Raisa~Jose, Noel~Alex~Samuel, Khadheeja~Safar, and Kalyanaraman~P%
\thanks{The authors are with the School of Computer Science and Engineering, Vellore Institute of Technology, Vellore, India.}}

\maketitle

\begin{abstract}
\vTwelveAbstract
\end{abstract}

\begin{IEEEkeywords}
code review, annotation anchoring, source code tracking, software evolution, language models
\end{IEEEkeywords}

\section{Introduction}
A review comment refers both to a location and to a concern. When code is revised, the location may move while the concern remains, or the concern may be addressed while most of the quoted code survives. In five open-source Python projects, \wpeOutdated{} of \wpeThreads{} inline review threads on merged pull requests (\wpeOutdatedPct\%; 95\% pull-request cluster interval \wpeOutdatedLo--\wpeOutdatedHi\%) were marked outdated by GitHub, and \wpeNullLinePct\% of the outdated threads had no position in the current diff, so the reviewer sees the comment but not where it now applies.

\sys{} is the resubmission model of a capstone review portal in which faculty supervisors review a team's code and report over several versions. Consider a supervisor who comments on version~1 that a retry loop in \texttt{client.py} never backs off. The team has an identical copy of the loop in \texttt{worker.py}. In version~2 it fixes the loop in \texttt{client.py} and leaves the copy alone. A quotation-based locator no longer finds the reviewed text in place, finds it unchanged in \texttt{worker.py} and reports the comment \lbl{moved}. The supervisor then sees the old comment on unchanged code, with nothing to say that the reviewed code changed, and may keep or close the concern without seeing the fix. This is the mechanism behind every missed re-check of the portal's deployed locator on real histories (Section~\ref{sec:layer}). GitHub would instead mark the thread outdated and drop its position. In the configuration evaluated here (the deployed portal still runs the cascade alone, Section~\ref{sec:impl}), \sys{} keeps the placement, but its copy guard notices that the quoted text was already in \texttt{worker.py} and asks for a re-check, shows the reviewed quotation beside the current code, and keeps the issue open until the supervisor decides.

The example separates three properties that tools usually merge: whether the reviewed code \emph{changed}, whether the anchor still points at it (\emph{validity}), and whether the supervisor's \emph{concern} is unresolved. The first two are properties of code and can be measured; the third is known only to the supervisor, and a change is evidence that a concern may be affected, not that it has been addressed \cite{anand2026impact}. Whether carrying feedback forward helps supervisors re-review and students act on feedback is motivated by feedback-literacy research \cite{carless2018feedback,winstone2017supporting} but not answered by it. This paper addresses the technical problem and the issue-state rules that connect it to supervision; it measures no learning, re-review time or workload outcome.

Content-bearing anchors, fuzzy matching, diff-based tracking and region mapping are established \cite{brush2001robust,w3c2017annotation,asaduzzaman2013lhdiff,falleri2014fine,hu2026codemapper}, and review tools already port or outdate comments \cite{gerritporting,github2018outdating}. We claim no new matching principle, and the portal's own matcher is not the best available. The contribution is \emph{safety-aware feedback persistence}: a contract under which relocation evidence may move an issue only towards a supervisor's attention, stated as five requirements, and a \emph{safety layer} that applies it around any locator. Relocation results become change labels computed from the reviewed quotation, two guards inspect the returned window, the reviewed quotation is kept, and no relocation resolves an issue. The layer meets S1--S3 under stated assumptions; the implemented lifecycle does not yet meet S4 and S5, and we specify and model-check a contract that does (Section~\ref{sec:statecheck}). Because the layer reads only the window a locator returns, it can be placed over a stronger region mapper or a language model.

Outcomes are ordered by consequence. A wrongly placed issue misdirects a supervisor, whereas a missing location is visible and prompts a search, so the primary safety failures are \emph{missed re-checks}, changed or deleted code not returned to the supervisor, and \emph{silent misplacements}, an issue shown on wrong code without a request; \emph{misdirected} requests, \emph{lost} anchors and \emph{unnecessary re-checks} cost effort. Every recommended configuration is reported with these and its placement accuracy, six outcomes with their own denominators; acceptable rates are an institutional choice, so we report trade-offs rather than a pass mark.

\textbf{Research questions.} \textbf{RQ1} (location): how accurately do the deployed cascade, conventional review policies, a contextual line tracker, a region mapper and language models re-find anchored code on controlled edits, real Python histories, student-authored repositories, a human-annotated benchmark and applied review suggestions in five languages? \textbf{RQ2} (safety layer): how often does carry-forward fail to return changed code or misplace an issue silently, by which mechanism, does a locator-independent layer reduce these failures over every locator and at what cost, and which configuration should a deployment choose for projects it has not seen? \textbf{RQ3} (state contract): does the implemented lifecycle meet S4 and S5, and what does a contract that meets them require?

\textbf{Contributions.} (1)~A persistence contract: the three properties, five safety requirements with their assumptions, a counterexample for every failure mechanism found (Section~\ref{sec:props}), and a comparison with review tools and trackers on each requirement (Section~\ref{sec:related}). (2)~A safety layer that applies the contract around any locator, built into a deployed portal whose system of record is still the cascade alone; the guards and identity rule are evaluated, not yet deployed (Sections~\ref{sec:design}--\ref{sec:impl}). (3)~An evaluation on five populations graded by evidential strength, including an oracle built from applied review suggestions, with all six outcomes for every configuration, ablations of the guards, the identity floor and copy placement, simulated supervisor interventions, and configuration selection on held-out projects (Sections~\ref{sec:method}--\ref{sec:layer}). Its negative results: a language model's own judgment of change is not a safe signal, a provenance-based copy guard fails where a copy guard is needed, the guards flag wrong placements rather than correct them, and insertions around an unchanged quotation evade every locator. (4)~A check of the implemented lifecycle by static scan, bounded model and scenario tests, and a stronger contract that closes the gaps it found (Section~\ref{sec:statecheck}).

An extended technical report in the artifact gives every analysis summarized here.
